\section{Experiments}
\label{sec:experiments}

\paragraph{Setup.}
The main experiments use Qwen2.5-VL-7B-Instruct \citep{qwen25vl2025}. We use ChartQA validation (1{,}920 examples) for development and held-out ChartQA test (2{,}500 examples) for the main task result \citep{masry2022chartqa}. We also evaluate InternVL2.5-1B/8B/26B \citep{internvl25} and numeric subsets of PlotQA \citep{methani2020plotqa}, ChartBench \citep{chartbench2024}, MathVista \citep{lu2023mathvista}, and InfographicVQA \citep{mathew2022infographicvqa}. Numeric answers are scored with the repository evaluator: normalized exact match for text answers and 5\% tolerance for numeric answers. The final Qwen scale adapters use $K=300$ pairs, 300 steps, learning rate $10^{-6}$, $\beta=2.0$, and margin $m=0.10$. All accuracy differences are in percentage points.

\subsection{Held-Out ChartQA and Verified Strict Errors}

The first question is deliberately narrow: can a scale-directed update fix natural scale errors without damaging the main task, and is the effect due to the scale objective rather than extra training?

\begin{table}[t]
\centering
\small
\resizebox{\linewidth}{!}{\begin{tabular}{llrrr}
\toprule
System & Starting point & Test acc. & Std. & Delta \\
\midrule
Raw Qwen2.5-VL-7B & Raw backbone & 0.7804 & -- & -- \\
Answer-tuned LoRA & Answer-only & 0.8860 & -- & +0.1056 \\
Answer-tuned + large-positive scale & Answer-tuned & 0.8932 & 0.0003 & +0.0072 \\
Clean direct large-positive scale & Raw backbone & 0.7456 & -- & -0.0348 \\
Clean direct small-positive scale & Raw backbone & 0.7136 & -- & -0.0668 \\
\bottomrule
\end{tabular}
}
\caption{Held-out ChartQA test attribution for Qwen2.5-VL-7B. The scale update is measured on top of the answer-tuned adapter; clean direct scale-only tuning is a safety control. ``Large-positive'' denotes the prefer-large adapter.}
\label{tab:official_test}
\end{table}

Table~\ref{tab:official_test} separates answer tuning from scale tuning. The raw pipeline obtains 78.04 on ChartQA test; answer tuning reaches 88.60. Continuing from that adapter, prefer-large scale-contrastive tuning reaches 89.32 over five seeds, +0.72 points with paired bootstrap 95\% CI $[+0.38,+1.09]$. Clean scale-only adaptation from the raw backbone hurts test accuracy by 3.5--6.7 points. The aggregate gain is small because strict scale errors are rare (Table~\ref{tab:strict_prevalence}); non-scale accuracy moves only from 89.35 to 89.66 (Table~\ref{tab:test_subset_results}), and TextVQA-500 is unchanged at 94.20 (a small adjacent-task check, not a general safety guarantee).

\begin{table}[t]
\centering
\small
\resizebox{\linewidth}{!}{\begin{tabular}{lcc}
\toprule
Continuation from answer-tuned adapter (300 steps, 3 seeds) & ChartQA val.\ acc. & Verified strict test errors recovered \\
\midrule
None (answer-tuned baseline) & \verify{0.8640} & \verify{0/17} \; [0.000, 0.184] \\
Ordinary answer replay & \verify{$\approx$0.8640} & \verify{1/17} \; [0.010, 0.270] \\
Random / shuffled numeric negatives & \verify{$\approx$0.8640} & \verify{[fill]} \\
Non-directional balanced pairs & \verify{$\approx$0.8640} & \verify{[fill]} \\
Directional scale pairs (prefer-large) & \verify{0.8695} & \verify{12/17} \; [0.469, 0.867] \\
Prefer-large adapter, main setting (5 seeds) & -- & \verify{14/17} \; [0.590, 0.938] \\
\bottomrule
\end{tabular}
}
\caption{Compute-matched controls. Every row continues the same answer-tuned adapter for the same number of steps; only the continuation data differs. Recovery is reported as counts over the \verify{17} blind-verified strict test errors with Wilson 95\% intervals.}
\label{tab:controls}
\end{table}

Table~\ref{tab:controls} answers the attribution question. On the verified strict test errors the answer-tuned baseline recovers \verify{0/17}. Continuing with ordinary answer examples recovers \verify{1/17}, and non-directional or random-negative continuation is statistically tied with the baseline on ChartQA validation. Directional scale supervision recovers \verify{12/17} in the three-seed control setting and \verify{14/17} with the main five-seed prefer-large adapter. The intervals are wide, but they exclude the control outcomes: the targeted effect comes from directional scale supervision, not from 300 more steps. This is a verified diagnostic-slice result, not a population-level rate (per-seed counts: Appendix~\ref{app:strict_verified}).

\subsection{ScaleQA-Natural}
\label{sec:scaleqa_results}

The second question is whether the same intervention is safe when scale conflicts are collected across regimes rather than filtered to one verified recovery subset.

\begin{table*}[t]
\centering
\small
\resizebox{\textwidth}{!}{\begin{tabular}{lrrrrrrrrr}
\toprule
Set & N & Raw & Ans. & Prompt & Large & Bal. & All & Real & Small \\
\midrule
ScaleQA-Natural-520 & 473 & 0.4567 & 0.5095 & 0.4419 & 0.4951 & -- & -- & -- & -- \\
ScaleQA-Core & 403 & 0.4591 & 0.4963 & 0.4293 & 0.4844 & 0.5092 & 0.5047 & 0.5002 & 0.4174 \\
InternVL-mined & 192 & 0.1510 & 0.3281 & 0.0990 & 0.1896 & 0.3323 & 0.2771 & 0.2510 & 0.7490 \\
\bottomrule
\end{tabular}
}
\caption{ScaleQA-Natural diagnostic results. The first row is the image-evaluable ScaleQA-Natural-473 subset; ScaleQA-Core and InternVL-mined report five-seed pair-mode means. Large/Small denote prefer-large/prefer-small adapters.}
\label{tab:scaleqa_results}
\end{table*}

On ScaleQA-Natural-473, answer tuning improves the raw backbone from 45.67 to 50.95, a scale-aware prompt drops to 44.19, and the global prefer-large adapter reaches 49.51, \emph{below} the answer-tuned base. On ScaleQA-Core, balanced and all-pair variants outperform the answer-tuned base; on the InternVL-mined subset, prefer-small adaptation is far stronger than every other global policy (74.90 versus 33.23 balanced).

\begin{table*}[t]
\centering
\small
\resizebox{\textwidth}{!}{\begin{tabular}{lrrrrrrr}
\toprule
ScaleQA-Natural-473 subset & $N$ & Answer-tuned & Prefer-large & Self-check & $\Delta$(large$-$ans.) [95\% CI] & Fix & Break \\
\midrule
Items from ChartQA training rows & 363 & 50.96 & 46.89 & 54.27 & $-4.1$ [$-7.8$, $-0.4$] & 4.7 & 8.8 \\
Items from ChartQA val/test rows (row-disjoint) & 110 & 50.91 & 58.18 & 55.45 & $+7.3$ [$+1.5$, $+13.6$] & 9.1 & 2.7 \\
All & 473 & 50.95 & 49.51 & 54.55 & $-1.4$ [$-4.7$, $+1.7$] & 5.7 & 7.4 \\
\midrule
\multicolumn{8}{l}{\emph{By direction} (answer-tuned / prefer-large / self-check):} \\
Prediction too large, training rows & 205 & 30.7 & 15.3 & 45.4 & & & \\
Prediction too small, training rows & 158 & 77.2 & 87.8 & 65.8 & & & \\
Prediction too large, row-disjoint & 36 & 22.2 & 14.4 & 33.3 & & & \\
Prediction too small, row-disjoint & 74 & 64.9 & 79.5 & 66.2 & & & \\
\bottomrule
\end{tabular}
}
\caption{Row-overlap analysis on ScaleQA-Natural-473. Prefer-large is the five-seed mean; $\Delta$ is the paired per-item difference with a bootstrap 95\% CI; Fix/Break are the percentages of items the adapter corrects/breaks relative to the answer-tuned model. The direction asymmetry is the same on both subsets; the aggregate sign differs because 56\% of training-row items are prediction-too-large versus 33\% of row-disjoint items.}
\label{tab:row_overlap}
\end{table*}

Table~\ref{tab:row_overlap} separates the items by whether their ChartQA row was seen during answer tuning. The aggregate ``adapter hurts'' result is carried entirely by the training-row items ($-4.1$ points, CI $[-7.8,-0.4]$, breaking 8.8\%). On the 110 row-disjoint items the same adapter \emph{helps} ($+7.3$ points, CI $[+1.5,+13.6]$, fixing 9.1\% and breaking 2.7\%) and edges self-check (58.18 versus 55.45). The direction breakdown shows why: in both subsets prefer-large helps when the prediction is too small and hurts when it is already too large; the subsets differ mainly in their mix of directions. Directionality is the stable finding; the sign of a global adapter's aggregate effect on a mined set depends on the set's direction composition, itself a reason not to deploy unconditional correction.

Training-free prompting baselines (Table~\ref{tab:tta_baselines_app}, appendix) are stronger than the global adapter on the full set. A self-check prompt reaches 54.55; majority voting adds little, while oracle prompt choice reaches 67.44. The question becomes whether the right answer is missing or merely not selected.

\begin{table}[t]
\centering
\small
\resizebox{\linewidth}{!}{\begin{tabular}{lc}
\toprule
Candidate source on ScaleQA-Natural-473 & Accuracy / coverage \\
\midrule
Single greedy sample (answer-tuned) & \verify{0.524} \\
Blind pass@20, no scale prompt, no transforms & \verify{0.793} \\
Prompt + adapter candidates (policy-pool oracle) & 0.9004 \\
Top-8 candidates incl.\ deterministic scale variants & 0.9175 \\
Augmented structured oracle & 0.9260 \\
\bottomrule
\end{tabular}
}
\caption{Candidate coverage on ScaleQA-Natural-473 decomposed by candidate source. Blind sampling uses no scale prompt and no deterministic transformation, so its coverage is not implied by the mining rule.}
\label{tab:blind_oracle}
\end{table}

Table~\ref{tab:blind_oracle} addresses a construction concern directly: because ScaleQA-Natural is mined through multiplicative relations and the structured oracle pools contain explicit scale variants, oracles over those pools are high partly by construction. The blind-generation row removes both. With no scale prompt and no transformation, \verify{20} stochastic samples already contain a correct answer for \verify{79.3\%} of examples, against \verify{52.4\%} for a single sample. Deterministic variants raise coverage only further, to 91.75--92.60; most of the gap to deployable methods is present without them.

\subsection{Cross-Backbone and External Checks}

The final question is scope: does the phenomenon extend beyond the main Qwen setting, and does the intervention transfer with it? Table~\ref{tab:transfer_boundary} (appendix) summarizes the picture. InternVL2.5-1B and 8B can be adapted with the same scale-pair recipe, and InternVL2.5-26B serves as a stronger diagnostic backbone on ScaleQA-Natural, PlotQA, ChartBench, MathVista, and InfographicVQA subsets. These checks establish two narrow facts: ratio-related scale errors occur under every backbone we test, and the prefer-large adapter does not transfer as a correction (InternVL2.5-8B shows only tiny differences among directions, and Qwen external subsets are unchanged across adapters). Because models, prompts, and subset sizes differ, they do not establish a common prevalence or taxonomy across models; we use them to bound the method's scope.
